\documentclass[12pt]{article}
\usepackage{xeCJK}
\usepackage{fontspec}
\usepackage{graphicx}
\usepackage{amsmath}
\usepackage{geometry}
\usepackage{fancyhdr}
\usepackage{indentfirst}
\usepackage{caption}
\usepackage{tabularx, booktabs}
\usepackage{amssymb}
\usepackage{amsfonts}
\geometry{a4paper, margin=1in}
\setCJKmainfont{Droid Sans Fallback}
\setmainfont{DejaVu Serif}
\setsansfont{Lato}
\setmonofont{Latin Modern Mono}
\pagestyle{fancy}
\setlength{\parindent}{2em}
\begin{document}

% ==== page 0 header/footer ====
\fancyhead[L]{Erstesein.eitendesBuch }
\fancyfoot[C]{}

% [Unknown block: number] 6



\par des geistigen Lebens nicht von der psycho-physischen Lebenseinheit der Menschennatur getrennt. Eine Theorie, welche die gesellschaftlich-geschichtlichen Tatsachen beschreiben und analysieren will,kann nicht von dieser Totalitat der Menschennatur absehen und sich auf das Geistige einschranken. Aber der Ausdruck teilt diesen Mangel mit jedem anderen,der angewandtworden ist; Gesellschaftswissenschaft (Soziologie),moralische,geschichtliche,Kulturwissenschaften:allediese Bezeichnungen leiden an demselben Fehler,zu eng zu sein in bezug auf den Gegenstand, den sie ausdrücken sollen. Und der hier gewahlte Name hat wenigstens den Vorzug, den zentralen Tatsachenkreis angemessen zu bezeichnen, von welchem aus in Wirklichkeit die Einheit dieser Wissenschaften gesehen, ihr Umfang entworfen, ihre Abgrenzung gegen die Naturwissenschaften, wenn auch noch so unvollkommen,vollzogenworden ist.


\par Der Beweggrund namlich, von welchem die Gewohnheit ausgegangen ist,diese Wissenschaften als eine Einheit von denen derNatur abzugrenzenreicht in die Tiefe und Totalitat des menschlichen SelbstbewuBtseins. Unangerührt noch von Untersuchungen iber den Ursprung des Geistigen,findet der Mensch in diesem SelbstbewuBtsein eine Souveranitat des Willens,eine Verantwortlichkeit derHandlungen,ein Vermogen, alles dem Gedanken zu unterwerfen und allem innerhalb der Burgfreiheit seiner Person zu widerstehen,durch welche er sich von der ganzen Natur absondert. Erfindet sich in dieser Natur in der Tat,einen Ausdruck Spinozas zu gebrauchen, als imperium in imperio. Und da für ihnnur das besteht, was Tatsache seines BewuBtseins ist,so liegt in dieser selbstandig in ihm wirkenden geistigen Welt jeder Wert,jeder Zweck des Lebens,in der Herstellung geistiger Tatbestande jedes Ziel seiner Handlungen. So sondert er yon dem Reich der Natur ein Reich der Geschichte,in welchem, mitten in demZusammenhang einer objektiven Notwendigkeit, welcher Natur ist,Freiheit an unzahligen Punkten dieses Ganzen aufblitzt;hier bringen die Taten des Willens, im Gegensatz zu dem mechanischen Ablauf der Naturveranderungen, welcher im Ansatz alles, was inihm erfolgt,schon enthalt, durch ihren Kraftaufwand und ihre Opfer,deren Bedeutung das Individuum ja in seiner Erfahrung gegenwirtig besitzt,wirklich etwas hervor, erarbeiten Entwicklung,in der Person undin der Menschheit: über die leere und ode Wiederholung von Naturlauf


% [Unknown block: footnote] Sehr genial drückt Pascal dies Lebensgefhlaus: Pensees Art.I.Toutes ces miseres—prouvent sa grandeur. Ce sontmiseres de grand seigneur,miseresdnroi depossede. (3) Nous avons une si grande idee de lame de lhomme, que nous ne pouvons souffrir d'en etre meprises,et den'etre pas dans l'estime d'une ame(s)(OeuvresParis 1866I,248,249).



\clearpage

\end{document}